
\section{Sistemas matemáticos comparativos: Vaidya}

El espaciotiempo de Vaidya entrante puede escribirse

\begin{equation}
ds^2
=
-
\left(
1-\frac{2m(v)}{r}
\right)
dv^2
+
2\,dv\,dr
+
r^2d\Omega^2
\cite{vaidya1951}.
\end{equation}

Las geodésicas radiales nulas salientes satisfacen

\begin{equation}
\boxed{
\frac{dr}{dv}
=
\frac12
\left(
1-\frac{2m(v)}{r}
\right).
}
\end{equation}

Una frontera aparente se localiza en

\begin{equation}
r_{\mathrm{AH}}=2m(v).
\end{equation}

La transferencia a IRL es exclusivamente metodológica: información global puede restringir accesibilidad en estados anteriores. Vaidya no constituye evidencia arqueológica.

\section{Sistema comparativo: Kerr}

Para geodésicas nulas ecuatoriales, \(E=1\), \(b=L/E\) y \(Q=0\),

\begin{equation}
\Delta=r^2-2Mr+a^2,
\end{equation}

\begin{equation}
\boxed{
R(r;b)
=
[r^2+a^2-ab]^2
-
\Delta(b-a)^2.
}
\end{equation}

El movimiento radial permitido satisface

\begin{equation}
R(r;b)\ge0.
\end{equation}

Una órbita nula circular crítica cumple

\begin{equation}
\boxed{
R=0,
\qquad
\partial_rR=0.
}
\end{equation}

Al cruzar la separatriz puede cambiar la conectividad del conjunto radial permitido \cite{kerr1963,wald1984}.

La analogía útil es la topología de accesibilidad, no una equivalencia física entre arquitectura y agujeros negros.

\section{Objeto definido por complemento}

Generalizamos mediante una relación \(\mathcal R\):

\begin{definition}[Objeto por complemento]
\begin{equation}
\boxed{
N_{\mathcal R}(\lambda)
=
\Omega
\setminus
\mathrm{Accessible}_{\mathcal R}(\Omega,\lambda).
}
\end{equation}
\end{definition}

\(\mathcal R\) puede significar accesibilidad morfológica, acceso de una pieza en \(\SE(3)\), escape causal o región permitida en espacio de fases.

Para

\begin{equation}
N_\lambda
=
\{x:\Phi(x;\lambda)<0\},
\end{equation}

un candidato crítico satisface

\begin{equation}
\boxed{
\Phi(x^\star;\lambda^\star)=0,
\qquad
\nabla_x\Phi(x^\star;\lambda^\star)=0,
}
\end{equation}

con condiciones de no degeneración específicas de cada sistema.
